\documentclass[10pt,a4paper]{amsart}
\usepackage[margin=19mm]{geometry}
\usepackage{amsmath,amssymb,mathtools}
\usepackage[scr=rsfs,cal=euler]{mathalfa}
\usepackage{xfrac}
\usepackage[unicode,hidelinks,bookmarks=false]{hyperref}
\newcommand{\ie}{i.e.\ }
\newcommand{\cf}{cf.\ }
\newcommand{\resp}{respectively\ }
\newcommand{\Subsection}{\subsection}
\providecommand{\nobreakdash}{\nobreak\hbox{-}}
\setlength{\emergencystretch}{2em}
\multlinegap0pt
\usepackage{bm}
\usepackage{bbm}
\usepackage{dsfont}
\usepackage{xspace}
\usepackage{cancel}
\usepackage{mathtools}
\usepackage{amsthm}
\makeatletter
\@ifundefined{Theorem}{\newtheorem{Theorem}{Theorem}}{}
\@ifundefined{Lemma}{\newtheorem{Lemma}[Theorem]{Lemma}}{}
\makeatother
\def\abs#1{\mathopen|#1\mathclose|}
\newcommand{\INTO}{\hookrightarrow} % embedding
\newcommand{\Ll}{{\mathcal{L}}} % Lagrangian planes
\newcommand{\N}{{\mathbb{N}}}
\newcommand{\R}{{\mathbb{R}}}
\newcommand{\norm}{{\rm norm}}
\newcommand{\p}{{\partial}}
\renewcommand{\SS}{{\mathbb{S}}}
\title[The Linearized Floer Equation in a Chart]{The Linearized Floer Equation in a Chart}
\author{Urs Frauenfelder, Joa Weber}
\date{Source: arXiv:2510.20945v2 (CC BY 4.0)}
\begin{document}
\maketitle

\noindent\emph{Source and licence.} The Linearized Floer Equation in a Chart. Version arXiv:2510.20945v2 (CC BY 4.0), distributed under the Creative Commons Attribution 4.0 International licence, as recorded in the arXiv licence metadata for this exact version and shown on the abstract page https://arxiv.org/abs/2510.20945v2. Full rights notice: licenses/arxiv-2510.20945-CC-BY-4.0.txt.\par\smallskip

\noindent\textit{Editorial scope.} Counted unit: the Theorem whose source label is \texttt{thm:paradarboux-Hess} in the article (source0.tex lines 2889--2898, source label \texttt{thm:paradarboux-Hess}), delivered together with the article's own complete original proof of it (source0.tex lines 2900--3126, one \texttt{proof} environment of 227 lines). No mathematics is written, repaired or re-derived: statement and proof are the article's own text line for line. Required context retained uncounted: eq:hyp-F (lines 2212--2214); eq:(C) (lines 2233--2237); eq:Behz-Hol-2 (lines 2842--2847); eq:hjghjghj77799 (lines 2871--2874); le:sum-quad (lines 3130--3136); thm:Fredholm-F\_C (lines 4097--4109). Every definition, display and label of those lines is retained unchanged. Omitted from this package and not redistributed: 1--2211, 2215--2232, 2238--2841, 2848--2870, 2875--2888, 2899--2899, 3127--3129, 3137--4096, 4110--4947. Nothing inside a retained excerpt is omitted or altered: every retained body is delivered exactly as the article's lines are written, so no omission receipt of any kind accompanies this package. Citations: the retained excerpts carry no citation command at all, so no bibliography entry of the article is transcribed and no reference is redistributed. Reference adaptation: none was needed and none was made; every reference inside the delivered unit resolves inside it, so no unresolved reference or citation is produced. Printed slips kept literal: no printed word, symbol, hypothesis or number of the counted statement or of its proof was altered. Layout and class: the article's own class is \texttt{sigma} with packages including ifpdf, scrextend, hycolor, xcolor; the wrapper is the lane's pinned \texttt{amsart} editorial wrapper at 10pt on A4, which restates the theorem-like environments the retained text needs and adds the article's own macro definitions that the retained text needs (\texttt{\textbackslash INTO}, \texttt{\textbackslash Ll}, \texttt{\textbackslash N}, \texttt{\textbackslash R}, \texttt{\textbackslash SS}, \texttt{\textbackslash abs}, \texttt{\textbackslash norm}, \texttt{\textbackslash p}), with extra wrapper packages (bm, bbm, dsfont, xspace, cancel, mathtools). The delivered numbering is the unit's own. Assets: no retained span contains a figure environment, a graphics-inclusion command, a PDF-page inclusion, a table or a TikZ picture; no image file is copied into this package, the package declares no asset dependency and no image file is redistributed with it. Licence: version arXiv:2510.20945v2 of this article is distributed under the Creative Commons Attribution 4.0 International licence, as recorded in the arXiv licence metadata for this exact version and shown on the abstract page https://arxiv.org/abs/2510.20945v2. A full rights notice is delivered at licenses/arxiv-2510.20945-CC-BY-4.0.txt. Characters: every retained excerpt is pure ASCII, so the delivered engine, XeLaTeX with its default Latin Modern text font, reports no missing character.
\begin{equation}\label{eq:hyp-F}
 F\in C^0(U_1,\Ll(H_1,H_0)\cap\Ll(H_2,H_1))
\end{equation}

 \begin{equation}\itemsep=0pt\label{eq:(C)}
 \norm{C^v-C^w}_{\Ll(H_1)}
 \le\kappa\bigl(\abs{v-w}_{H_2}
 +\min\{\abs{v}_{H_2},\abs{w}_{H_2}\}\cdot\abs{v-w}_{H_1}\bigr).
 \end{equation}

\begin{equation}\label{eq:Behz-Hol-2}
 W^{r,2}\times L^2\stackrel{r>\frac12}{\longrightarrow} L^2
 ,\qquad
 W^{1,2}\times W^{1,2}\to W^{1,2}
 .
\end{equation}

\begin{equation}\label{eq:hjghjghj77799}
 B_u^{-1}\p_t\colon\ H_1\to H_0,\qquad H_2\to H_1
 ,\qquad \forall u\in U_1
\end{equation}

\begin{Lemma}\label{le:sum-quad}
For $k\in\N$ real numbers $a_1,\dots,a_k>0$, there is the inequality
\[
 \Biggl(\sum_{j=1}^k a_j\Biggr)^2
 \le k\sum_{j=1}^k a_j^2 .
\]
\end{Lemma}

\begin{Theorem}\label{thm:Fredholm-F_C}
Fix a loop of symplectic bilinear forms
\smash{$\omega\in W^{1,2}\bigl(\SS^1,\Omega\bigl(\R^{2n}\bigr)\bigr)$}.
Then the operator
\[
 F_{B_\omega^{-1}}=B_\omega^{-1}\p_t
 \colon\ H_{k}\to H_{k-1}
 ,\qquad
 H_k:=W^{k,2}\bigl(\SS^1,\R^{2n}\bigr) ,
\]
is Fredholm of index zero
whenever $k=1,2$.
\end{Theorem}

\begin{Theorem}\label{thm:paradarboux-Hess}
The para-Darboux weak Hessian field $A_0$ is almost extendable.
Moreover, the pair $(F,C)$ defined for $u\in U_1$ by
\[
 F^u\colon\ \eta\mapsto B_u^{-1}\dot \eta
 ,\qquad
 C^u\colon\ \eta\mapsto \bar L_u(\eta,\dot u) ,
\]
is a decomposition.
\end{Theorem}

\begin{proof}[Proof of Theorem~\ref{thm:paradarboux-Hess}]
The proof has four Steps~1, 2, (\texttt{C}), and (\texttt{F}).

\medskip\noindent
\textbf{Step~1.}
By~(\ref{eq:Behz-Hol-2}),
the map $u\mapsto C^u$ is element of the space
$C^0(U_1,\Ll(\underline{H_r},H_0))$ and of the space
$C^0(U_2,\Ll(\underline{H_1},H_1))$,
where $H_r=W^{r,2}\bigl(\SS^1,\R^{2n}\bigr)$.


\medskip\noindent
\textbf{Step~2.}
We need to show that the map $F\colon u\mapsto F^u=B_u^{-1}\p_t$
is element of the space $C^0(U_1,\Ll(H_1,H_0)\cap\Ll(H_2,H_1))$;
see~(\ref{eq:hyp-F}).

\begin{proof}
We prove that $F\in C^0(U_1,\Ll(H_2,H_1))$;
similarly $F\in C^0(U_1,\Ll(H_1,H_0))$.

Fix $u\in U_1$ and let $K=K(u)$ be a compact neighborhood
of the (compact) image of $u$ in~${\mathfrak{U}\subset\R^{2n}}$.
Pick $v\in U_1$ near $u$ taking values in $K$ as well.
The smooth map $x\mapsto B_x^{-1}$ is Lipschitz continuous
since $K=K(u)$ is compact. Let $\lambda_u$ be the Lipschitz constant.
Now, given any~${\xi\in H_2}$, we estimate
\begin{equation*}
\begin{split}
 \abs{(F^u-F^v)\xi}_{W^{1,2}}^2
 &=\bigl|\bigl(B_u^{-1}-B_v^{-1}\bigr)\dot\xi\bigr|_{L^2}^2
 +\bigl|
 \bigl(B_u^{-1}-B_v^{-1}\bigr)\ddot\xi
 +\bigl(\p_t\bigl(B_u^{-1}-B_v^{-1}\bigr)\bigr)\dot\xi
 \bigr|_{L^2}^2
\\
 &\le\lambda_u^2\abs{u-v}_{L^\infty}^2\abs{\xi}_{W^{1,2}}^2
 +2\lambda_u^2\abs{u-v}_{L^\infty}^2\abs{\xi}_{W^{2,2}}^2
 \\
 &\quad{}
 +4 \bigl|{\rm d}B_u^{-1}\bigl(\dot u\,\underline{-\,\dot v},\dot\xi\bigr)\bigr|_{L^2}^2
 +4 \bigl|\bigl(\underline{{\rm d}B_u^{-1}}-{\rm d}B_v^{-1}\bigr)\bigl(\dot v,\dot\xi\bigr)\bigr|_{L^2}^2
\\
 &\le3\lambda_u^2\abs{u-v}_{W^{1,2}}^2\abs{\xi}_{W^{2,2}}^2
 +4\bigl|{\rm d}B_u^{-1}\bigr|_{L^\infty(K)}^2\abs{u-v}_{W^{1,2}}^2\bigl|\dot\xi\bigr|_{L^\infty}^2
 \\
 &\quad{}
 +4\lambda_u^2\abs{u-v}_{L^\infty}^2\abs{v}_{W^{1,2}}^2\bigl|\dot\xi\bigr|_{L^\infty}^2
\\
 &\le\bigl(3\lambda_u^2+4\bigl|{\rm d}B_u^{-1}\bigr|_{L^\infty(K)}^2
 +4\lambda_u^2 \abs{v}_{W^{1,2}}^2\bigr)
 \abs{u-v}_{W^{1,2}}^2 \abs{\xi}_{W^{2,2}}^2.
\end{split}
\end{equation*}
Here inequality one uses Lipschitz continuity of $K\ni x\mapsto B_x^{-1}$
with Lipschitz constant $\lambda_u$. We also added \textit{zero}
and used the triangle inequality.
Inequality two uses the Sobolev embedding $W^{1,2}\INTO C^0$ with
constant $1$, we also used Lipschitz continuity again.
This proves Step~2.
\end{proof}


\noindent
\textbf{Step~(\texttt{C}).} The map $u\mapsto C^u$ satisfies the scale
Lipschitz estimate~(\ref{eq:(C)}).


\begin{proof}
It remains to check the local scale Lipschitz axiom~(\texttt{C}),
see~(\ref{eq:(C)}), for the map~$C$.
To this end suppose $u\in U_1=W^{1,2}\bigl(\SS^1,\mathfrak{U}\bigr)$.
Since $W^{1,2}\subset C^0$, the image of $u$ is a compact subset
of~$\mathfrak{U}$. Therefore, there exist, firstly, an open subset
$\mathfrak{V}$ which contains the image of $u$ and, secondly, a~constant $c$ such that
\begin{align}
 &\bigl\|\bar L_x\bigr\|_{\Ll(\R^{2n}\times\R^{2n};\R^{2n})}
\le c, \nonumber
\\
 &\bigl\|\bar L_x-\bar L_x\bigr\|_{\Ll(\R^{2n}\times\R^{2n};\R^{2n})}
\le c \abs{x-y}_{\R^{2n}}, \nonumber
\\
 &\bigl\| {\rm d}\bar L_x\bigr\|_{\Ll(\R^{2n}\times\R^{2n}\times\R^{2n};\R^{2n})}
\le c ,\nonumber
\\
 &\bigl\|{\rm d}\bar L_x-{\rm d}\bar L_x\bigr\|_{\Ll(\R^{2n}\times\R^{2n}\times\R^{2n};\R^{2n})}
\le c \abs{x-y}_{\R^{2n}},\label{eq:L-est}
\end{align}
whenever $x,y\in\mathfrak{V}$.
We define an open neighborhood of $u$ in $U_1$ by
\[
 V_u:=W^{1,2}\bigl(\SS^1,\mathfrak{V}\bigr)
 \cap B_1^{H_1}(u),
\]
where \smash{$B_1^{H_1}(u)$} is the open radius $1$ ball in $H_1$ centered at $u$.
Pick elements $v,w\in V_u\cap H_2 =V_u\cap W^{2,2}$.
Note that
\[
 \abs{v}_{W^{1,2}}\le \abs{u}_{W^{1,2}}+1
 ,\qquad
 \abs{w}_{W^{1,2}}\le \abs{u}_{W^{1,2}}+1 .
\]
Now for $\xi\in H_1$, we estimate
\begin{equation*}
\begin{split}
 \abs{(C^v-C^w)\xi}_{W^{1,2}}^2
 \le{}& \bigl|\bar L_v(\xi,\dot v)-\bar L_w(\xi,\dot w)\bigr|_{L^2}^2
 +\bigl|\p_t\bigl(\bar L_v(\xi,\dot v)-\bar L_w(\xi,\dot w)\bigr)\bigr|_{L^2}^2
\\
 \le{}& 2\bigl|\bar L_v(\xi,\dot v-\dot w)\bigr|_{L^2}^2
 +2\bigl|\bigl(\bar L_v-\bar L_w\bigr) (\xi,\dot w)\bigr|_{L^2}^2
 \\
 &{+}\,7\bigl|{\rm d}\bar L_v(\dot v-\dot w,\xi,\dot v)\bigr|_{L^2}^2
 +7\bigl|{\rm d}\bar L_v(\dot w,\xi,\dot v-\dot w)\bigr|_{L^2}^2
 \\
 &{+}\,7\bigl|\bigl({\rm d}\bar L_v-{\rm d}\bar L_w\bigr)(\dot w,\xi,\dot w)\bigr|_{L^2}^2
 \\
 &{+}\,7\bigl|\bar L_v\bigl(\dot \xi,\dot v-\dot w\bigr)\bigr|_{L^2}^2
 +7\bigl|\bigl(\bar L_v-\bar L_w\bigr)\bigl(\dot \xi,\dot w\bigr)\bigr|_{L^2}^2
 \\
 &{+}\,7\bigl|\bar L_v(\xi,\ddot v-\ddot w)\bigr|_{L^2}^2
 +7\bigl|\bigl(\bar L_v-\bar L_w\bigr)(\xi,\ddot w)\bigr|_{L^2}^2,
\end{split}
\end{equation*}
where in inequality two we added four times zero
and used Lemma~\ref{le:sum-quad} for $k=7$ summands.
Now we estimate summand by summand starting below
\begin{equation*}
\begin{split}
 \bigl|\bigl(\bar L_v-\bar L_w\bigr)(\xi,\ddot w)\bigr|_{L^2}^2
 &=\int_0^1\bigl|\bigl(\bar L_{v_t}-\bar L_{w_t}\bigr) (\xi_t,\ddot w_t)\bigr|^2\, {\rm d}t
\\
 &\le \int_0^1c^2
 \abs{v_t-w_t}^2\abs{\xi_t}^2\abs{\ddot w_t}^2\, {\rm d}t
\\
 &\le c^2 \abs{v-w}_{L^\infty}^2\abs{\xi}_{L^\infty}^2
 \abs{\ddot w}_{L^2}^2
\\
 &\le c^2 \abs{v-w}_{W^{1,2}}^2\abs{\xi}_{W^{1,2}}^2
 \abs{w}_{W^{2,2}}^2 .
\end{split}
\end{equation*}
Here inequality one uses~(\ref{eq:L-est}) and the Sobolev embedding
$L^\infty\INTO W^{1,2}$ with constant $1$.
Similarly, we estimate
\begin{equation*}
\begin{split}
%1
& \bigl|\bar L_v(\xi,\ddot v-\ddot w)\bigr|_{L^2}^2
\le c^2 \abs{\xi}_{W^{1,2}}^2\abs{v-w}_{W^{2,2}}^2,
\\
%2
& \bigl|\bigl(\bar L_v-\bar L_w\bigr)\bigl(\dot \xi,\dot w\bigr)\bigr|_{L^2}^2
\le c^2 \abs{v-w}_{W^{1,2}}^2 \abs{\xi}_{W^{1,2}}^2 \abs{w}_{W^{2,2}}^2,
\\
%3
& \bigl|\bar L_v\bigl(\dot \xi,\dot v-\dot w\bigr)\bigr|_{L^2}^2
\le c^2 \abs{\xi}_{W^{1,2}}^2 \abs{v-w}_{W^{2,2}}^2,
\\
%4
& \bigl|\bigl({\rm d}\bar L_v-{\rm d}\bar L_w\bigr)(\dot w,\xi,\dot w)\bigr|_{L^2}^2
\le c^2 \bigl(1+\abs{u}_{W^{1,2}} \bigr)^2
 \abs{v-w}_{W^{1,2}}^2
 \abs{\xi}_{W^{1,2}}^2\abs{w}_{W^{2,2}}^2,
\\
%5
& \bigl|{\rm d}\bar L_v(\dot w,\xi,\dot v-\dot w)\bigr|_{L^2}^2
\le c^2 \bigl(1+\abs{u}_{W^{1,2}} \bigr)^2 \abs{\xi}_{W^{1,2}}^2
 \abs{v-w}_{W^{2,2}}^2,
\\
%6
& \bigl|{\rm d}\bar L_v(\dot v-\dot w,\xi,\dot v)\bigr|_{L^2}^2
\le c^2 \bigl(1+\abs{u}_{W^{1,2}} \bigr)^2
 \abs{v-w}_{W^{2,2}}^2\abs{\xi}_{W^{1,2}}^2,
\\
%7
& \bigl|\bigl(\bar L_v-\bar L_w\bigr) (\xi,\dot w)\bigr|_{L^2}^2
\le c^2 \abs{v-w}_{W^{1,2}}^2 \abs{\xi}_{W^{1,2}}^2 \abs{w}_{W^{2,2}}^2,
\\
%8
& \bigl|\bar L_v(\xi,\dot v-\dot w)\bigr|_{L^2}^2
\le c^2 \abs{\xi}_{W^{1,2}}^2\abs{v-w}_{W^{2,2}}^2.
\end{split}
\end{equation*}
Continuing the above estimate
and returning to the notation $H_k=W^{k,2}$, we get
\begin{equation*}
\begin{split}
 \abs{(C^v-C^w)\xi}_{H_1}^2
 &\le \underbrace{7 c^2 \bigl(7+4\abs{u}_{H_1}^2\bigr)}_{=:\kappa^2}
 \bigl(\abs{v-w}_{H_2}^2+\abs{w}_{H_2}^2\abs{v-w}_{H_1}^2\bigr)
 \abs{\xi}_{H_1}^2 .
\end{split}
\end{equation*}
This implies that the operator norm is bounded by
\[
 \norm{C^v-C^w}_{\Ll(H_1)}
 \le\kappa (\abs{v-w}_{H_2}+\abs{w}_{H_2}\abs{v-w}_{H_1} ).
\]
Interchanging the roles of $v$ and $w$, we get
\[
 \norm{C^v-C^w}_{\Ll(H_1)}
 \le\kappa (\abs{v-w}_{H_2}+\abs{v}_{H_2}\abs{v-w}_{H_1} ).
\]
The above two estimates imply the scale Lipschitz estimate
\begin{equation}\label{eq:C-paraDarboux}
 \norm{C^v-C^w}_{\Ll(H_1)}
 \le\kappa (\abs{v-w}_{H_2}
 +\min\{\abs{v}_{H_2},\abs{w}_{H_2}\}\abs{v-w}_{H_1}
 ),
\end{equation}
which is precisely~(\ref{eq:(C)}) in axiom~(\texttt{C}).
This proves Step~(\texttt{C}).
\end{proof}

\noindent
\textbf{Step~(\texttt{F}).} $\forall u\in U_1\colon$
$F^u_2:=B_u^{-1}\p_t |_{H_2}\colon H_2\to H_1$ is Fredholm of index zero.

\begin{proof}
The proof was given in~(\ref{eq:hjghjghj77799})
as a consequence of Theorem~\ref{thm:Fredholm-F_C}.
\end{proof}

The proof of Theorem~\ref{thm:paradarboux-Hess} is complete.
\end{proof}

\end{document}
